\documentclass[a4paper]{arrowhead}

\usepackage[yyyymmdd]{datetime}
\usepackage{etoolbox}
\usepackage[utf8]{inputenc}
\usepackage[T1]{fontenc}
\usepackage{multirow}
\usepackage{longtable}
\usepackage{tikz}
\usetikzlibrary{positioning,shapes.multipart,arrows.meta}

\renewcommand{\dateseparator}{-}
\setlength{\emergencystretch}{3em}
\lstset{language=HTTP,breakatwhitespace=false,columns=fullflexible}

\newcommand{\fparam}[1]{\textit{\textcolor{ArrowheadBlue}{#1}}}

%% Special references
\newcommand{\mref}[2][]{{\textcolor{ArrowheadPurple}{\hyperref[sec:model:#2]{\ifstrempty{#1}{#2}{#1}}}}}
\newcommand{\oref}[1]{{\textcolor{ArrowheadBlue}{\hyperref[sec:operations:#1]{#1}}}}
\newcommand{\pdef}[1]{{\textcolor{ArrowheadGrey}{#1 \label{sec:model:primitives:#1} \label{sec:model:primitives:#1s}}}}
\newcommand{\pref}[1]{{\textcolor{ArrowheadGrey}{\hyperref[sec:model:primitives:#1]{#1}}}}

% as template_idd.tex; the input and output arguments are passed already formatted (\mref or \pref)
\newrobustcmd\osubsection[5]{
  \addtocounter{subsection}{1}
  \addcontentsline{toc}{subsection}{\protect\numberline{\thesubsection}operation \textcolor{ArrowheadBlue}{#1}}
  \subsection*{
    \thesubsection\quad
    #2 \textcolor{ArrowheadPurple}{#3} \\
    \small
    \hspace*{0.075\textwidth}\begin{minipage}{0.1\textwidth}
      \vspace*{1mm}
      Operation: \\
      \notblank{#4}{Input: \\}{}
      \notblank{#5}{Output: \\}{}
    \end{minipage}
    \begin{minipage}{0.825\textwidth}
      \vspace*{1mm}
      \textcolor{ArrowheadBlue}{#1} \\
      \notblank{#4}{#4 \\}{}
      \notblank{#5}{#5 \\}{}
    \end{minipage}
  }
  \label{sec:operations:#1}
}
\newrobustcmd\msubsection[2]{
  \addtocounter{subsection}{1}
  \addcontentsline{toc}{subsection}{\protect\numberline{\thesubsection}#1 \textcolor{ArrowheadPurple}{#2}}
  \subsection*{\thesubsection\quad#1 \textcolor{ArrowheadPurple}{#2}}
  \label{sec:model:#2} \label{sec:model:#2s}
}
%%

\usepackage{xurl}
\begin{document}
\sloppy

%% Arrowhead Document Properties
\ArrowheadTitle{sysMonitorData MQTT/TLS/JSON}
\ArrowheadServiceID{sysMonitorData}
\ArrowheadType{Interface Design Description}
\ArrowheadTypeShort{IDD}
\ArrowheadVersion{5.3.0}
\ArrowheadDate{2026-11-09}
\ArrowheadAuthor{Jerker Delsing}
\ArrowheadStatus{REVIEW DRAFT}
\ArrowheadContact{jerker.delsing@ltu.se}
\ArrowheadFooter{\href{www.arrowhead.eu}{www.arrowhead.eu}}
\ArrowheadSetup
%%

%% Front Page
\begin{center}
  \vspace*{1cm}
  \huge{\arrowtitle}

  \vspace*{0.2cm}
  \LARGE{\arrowtype}
  \vspace*{1cm}

  \Large{Service ID: \textit{"\arrowid"}}

  \vspace*{\fill}

  % Front Page Image

  \vspace*{1cm}
  \vspace*{\fill}

  % Front Page Abstract
  \begin{abstract}
    This document describes the MQTT protocol with TLS payload security and JSON payload encoding variant of the \texttt{sysMonitorData} service of the SysMonitor support system (SysMonitor-Data-MQTT-IDD-5.2.2).
  \end{abstract}

  \vspace*{1cm}
\end{center}

\newpage
%%

%% Table of Contents
\tableofcontents
\newpage
%%

\section{Documentation edition and implementation scope}
This is the 5.3.0 documentation edition, dated 2026-11-09. It carries forward the system and service contracts from the retained 5.2.2 documentation and identifies the current primary implementation sources. The former \path|documentation/v5.2.2| folder is retained unchanged.

Document version 5.3.0 is distinct from executable, dependency, registered system/service and wire-contract versions. Values still defined by the source (including 5.2.2 artefact versions, 1.0.0 service versions and \path|SysMonitor-AMQP-IDD-5.2.2| payload identifiers) are preserved. This documentation update does not change runtime behaviour or establish a new implementation release. Historical verification statements describe the tests recorded in the inherited baseline; this revision validates document compilation and references.

\section{Overview}
\label{sec:overview}

This document describes the \texttt{sysMonitorData} service interface, which provides the stream of all received observations, republished on the local MQTT broker.

It's implemented using protocol, encoding as stated in below table:
\begin{table}[ht!]
  \centering
  \begin{tabular}{|l|l|l|}
    \hline
    \rowcolor{gray!33} Profile type & Type & Version \\ \hline
    Transfer protocol & MQTT & 5.0 \\ \hline
    Data encryption & TLS & 1.3 \cite{rfc8446}, 1.2 \cite{rfc5246} \\ \hline
    Encoding & JSON & RFC 8259 \cite{rfc8259} \\ \hline
    Compression & N/A & - \\ \hline
    Semantics & N/A & - \\ \hline
    Ontology & N/A & - \\ \hline
  \end{tabular}
  \caption{Communication and semantics details used for the \texttt{sysMonitorData} service interface}
  \label{tab:comunication_semantics_profile}
\end{table}

This document provides the Interface Design Description IDD to the \textit{sysMonitorData -- Service Description} document \cite{smo-sysMonitorData-sd}.

\subsection{Interface registration}

\begin{lstlisting}[caption={Registered service instance (TestCloud).}]
{
  "serviceDefinitionName": "sysMonitorData",
  "version": "1.0.0",
  "metadata": {
    "contract": "SysMonitor-Data-MQTT-IDD-5.2.2",
    "schemaVersion": "1.0"
  },
  "interfaces": [
    {
      "templateName": "generic_mqtts",
      "protocol": "ssl",
      "policy": "CERT_AUTH",
      "properties": {
        "accessAddresses": [
          "mqtt-broker"
        ],
        "accessPort": 8883,
        "baseTopic": "arrowhead/TestCloud/sysmonitor/data",
        "operations": [
          "subscribe"
        ]
      }
    }
  ]
}
\end{lstlisting}

\section{Service Operations}
\label{sec:operations}

\osubsection{subscribe}{SUBSCRIBE}{\texttt{\textless{}base\allowbreak{}Topic\textgreater{}/\allowbreak{}\textless{}observation\allowbreak{}Type\textgreater{}/\allowbreak{}\textless{}provider\allowbreak{}System\textgreater{}}}{}{\mref{StoredObservation}}

A consumer connects to the MQTT broker with mutual TLS (its system certificate) and subscribes below the base topic, e.g. \texttt{arrowhead/\allowbreak{}Test\allowbreak{}Cloud/\allowbreak{}sysmonitor/\allowbreak{}data/\allowbreak{}echo/\allowbreak{}+} or \texttt{arrowhead/\allowbreak{}Test\allowbreak{}Cloud/\allowbreak{}sysmonitor/\allowbreak{}data/\allowbreak{}+/\allowbreak{}+}. SysMonitor publishes every stored observation with QoS 1, not retained, with \texttt{publishedTime} set; the broker's access list decides who may subscribe.

\begin{lstlisting}[caption={A published observation.}]
{
  "schemaVersion": "string",
  "observationId": "DataManager-echo-1792670400000-42",
  "providerSystem": "DataManager",
  "serviceInstance": "DataManager",
  "observationType": "echo",
  "sourceTime": "2026-10-22T12:00:00.000Z",
  "receivedTime": "2026-10-22T12:00:00.120Z",
  "publishedTime": "2026-10-22T12:00:00.300Z",
  "value": {
    "status": "alive"
  },
  "quality": "good",
  "pipeline": {
    "routingKey": "TestCloud.sysmonitor.DataManager.DataManager.echo"
  }
}
\end{lstlisting}

\newpage

\section{Data Models}
\label{sec:model}

\msubsection{struct}{StoredObservation}




\noindent\begin{tabularx}{\textwidth}{| >{\raggedright\arraybackslash}p{3.9cm} | >{\raggedright\arraybackslash\small}p{5.3cm} | X |} \hline
\rowcolor{gray!33} Object Field & Value Type & Description \\ \hline
"schema\allowbreak{}Version" & \pref{String} & \texttt{1.\allowbreak{}0}. \\ \hline
"observation\allowbreak{}Id" & \pref{String} & Identifier (publisher's, or a SHA-256 of routing key and content). \\ \hline
"provider\allowbreak{}System" & \pref{String} & Publishing system. \\ \hline
"service\allowbreak{}Instance" & \pref{String} & Service instance. \\ \hline
"observation\allowbreak{}Type" & \pref{String} & Observation type. \\ \hline
"source\allowbreak{}Time" & \pref{DateTime} & Publisher time, UTC. \\ \hline
"received\allowbreak{}Time" & \pref{DateTime} & Time SysMonitor received it. \\ \hline
"published\allowbreak{}Time" & \pref{DateTime} & Time it was forwarded on \texttt{sysMonitorData}. \\ \hline
"value" & \pref{Object} & Observed value. \\ \hline
"unit" & \pref{String} & Unit. \\ \hline
"quality" & \pref{String} & Quality. \\ \hline
"pipeline" & \pref{Object} & Intake details (routing key, exchange). \\ \hline
\end{tabularx}



\subsection{Primitives}
\label{sec:model:primitives}

\noindent\begin{tabularx}{\textwidth}{| p{4.6cm} | X |} \hline
\rowcolor{gray!33} Type & Description \\ \hline
\pdef{DateTime} & JSON String with an ISO 8601 UTC timestamp, e.g. \texttt{2026-10-22T12:00:00Z}. \\ \hline
\pdef{Number} & JSON Number, an integer. \\ \hline
\pdef{Object} & JSON Object; \pref{Object}$<$A$>$ has string keys and values of type A. \\ \hline
\pdef{String} & JSON String. \\ \hline
\end{tabularx}

\newpage

\section{Go and Rust implementation variants}
\subsection{Go implementation}
Primary system folder: \path|Eclipse-Arrowhead-go/ah5-SysMonitor-go|.\par
The system owns its Go module, launcher and build configuration. Use the per-system \path|Makefile| and \path|go.mod|; the containing Go tree is the build source, not a historical Arrowhead copy. Ordinary native roles use independent identities and JSON snapshot state rather than the Java database. Do not assume Java SQL migration, clustering or identical transaction/retry semantics. The registry has its own configuration and optional native AMQP monitor; other roles must be checked against their own registrations.
\paragraph{Listening port.} Built-in default 127.0.0.1:8475 (Arrowhead default). The shipped config.json, compose.yaml and Dockerfile use the same port on all interfaces ("listen": "0.0.0.0:8475"; 8540--8555 layout before 2026-11-09). Override with "listen" in the configuration file, SERVER\_PORT=N or -server.port N; the port overrides keep the configured host.
\paragraph{Native behavior and outstanding work.} Authorized AMQPS intake, strict observation validation, atomic dedup/history/outbox, subscriptions, history, MQTT QoS1 forwarding. Java MQTTv5 (Go uses 3.1.1), complete live-broker/UI parity, Java history import, broad schema golden tests. These are the boundaries recorded by the native compatibility document; later implementation changes still require direct qualification.
Native design/route evidence: \path|Eclipse-Arrowhead-go/ah5-SysMonitor-go/documentation/SysMonitor|.\par
For this Java interface, select only a native transport and operation explicitly implemented and advertised by that port. A matching method/path does not establish DTO, error, authorization, retry or persistence equivalence. Java MQTT/AMQP bindings are not automatically provided by a Go or Rust role. Differential contract and security checks remain necessary before substitution.
\subsection{Rust implementation}
Primary system folder: \path|Eclipse-Arrowhead-rust/ah5-SysMonitor-rust|.\par
The per-system entry point is declared as a binary of \path|arrowhead-support-rust| in \path|ah5-support-runtime-rust/Cargo.toml|. Build from the primary Rust workspace with the applicable package/binary selection. Documentation links can resolve into the shared core/support catalogue. Independent process identity and state remain required.\par
The ordinary native profile uses Tokio/Axum and rustls with peer-identity checks and independent locked JSON state. Network effects are not distributed database transactions. The precise trust, CRL reload and persistence guarantees are those of the native profile, not the Java implementation.\par
\paragraph{Listening port.} Default 127.0.0.1:8475 (Arrowhead default, role sysmonitor). Override with "listen": "host:port" in the configuration, SERVER\_PORT=N or --server.port N; the port overrides keep the configured host.
\paragraph{Native behavior and outstanding work.} Native mTLS AMQP intake, authorised binding reconciliation, revision subscriptions, SQLite history/outbox, confirmed MQTT output and retention. Full Java DTO/filter parity and wider operational qualification. These are the boundaries recorded by the native compatibility document; later implementation changes still require direct qualification.
Native design/route evidence: \path|Eclipse-Arrowhead-rust/ah5-support-runtime-rust/documentation/SysMonitor|.\par
For this Java interface, select only a native transport and operation explicitly implemented and advertised by that port. A matching method/path does not establish DTO, error, authorization, retry or persistence equivalence. Java MQTT/AMQP bindings are not automatically provided by a Go or Rust role. Differential contract and security checks remain necessary before substitution.

\bibliographystyle{IEEEtran}
\bibliography{bibliography}

\newpage

\section{Revision History}
\subsection{Amendments}

\noindent\begin{tabularx}{\textwidth}{| p{1cm} | p{3cm} | p{2cm} | X | p{4cm} |} \hline
\rowcolor{gray!33} No. & Date & Version & Subject of Amendments & Author \\ \hline

1 & 2026-10-22 & 5.2.2 & Document created from the Arrowhead LaTeX templates for the 5.2.2 Java Spring implementation, including the local extensions of this workspace. & Jerker Delsing \\ \hline
2 & 2026-10-30 & 5.2.2 & Checked against the primary Java source tree (Eclipse-Arrowhead-java): services, operations, paths and data models unchanged. & Claude Opus 5.5 (draft) \\ \hline
3 & 2026-11-09 & 5.3.0 & New documentation edition; primary-source references and Go/Rust implementation scope; 5.2.2 retained. & Documentation maintenance \\ \hline
\end{tabularx}

\subsection{Quality Assurance}

\noindent\begin{tabularx}{\textwidth}{| p{1cm} | p{3cm} | p{2cm} | X |} \hline
\rowcolor{gray!33} No. & Date & Version & Approved by \\ \hline

1 & & \arrowversion  &  \\ \hline

\end{tabularx}

\end{document}
